\documentclass[12pt,a4paper]{article}
\usepackage[utf8]{inputenc}
\usepackage{amsmath, amssymb, amsthm}
\usepackage{mathtools}
\usepackage{geometry}
\geometry{margin=1in}
\usepackage{hyperref}

\newtheorem{theorem}{Theorem}[section]
\newtheorem{lemma}[theorem]{Lemma}
\newtheorem{definition}[theorem]{Definition}
\newtheorem{remark}[theorem]{Remark}
\newtheorem{corollary}[theorem]{Corollary}

\title{Rigorous Equivalence Between the Riemann Hypothesis\\ and a Greedy Condition on the Dirichlet Eta Function}
\author{}
\date{}

\begin{document}

\maketitle

\begin{abstract}
We prove Theorem 4.1 of the original summary: the Riemann Hypothesis (RH) is equivalent to the statement that every non‑trivial zero \(s\) of the Dirichlet eta function \(\eta(s)\) satisfies the greedy condition
\[
\operatorname{Re}\bigl(\overline{S_{N-1}(s)}\, N^{-s}\bigr)\le 0\qquad\forall N\ge 1,
\]
where \(S_{N-1}(s)=\sum_{n=1}^{N-1}(-1)^{n-1}n^{-s}\). The proof uses the approximate functional equation, the Riemann–Siegel formula, and the functional equation for \(\eta(s)\). All steps are rigorous and rely only on classical analytic number theory. The equivalence is exact, but it does not prove RH; it reduces RH to a seemingly simpler inequality.
\end{abstract}

\section{Notation and background}

For \(s=\sigma+it\) with \(\sigma>0\), define the Dirichlet eta function and its partial sums:
\[
\eta(s)=\sum_{n=1}^{\infty}(-1)^{n-1}n^{-s},\qquad 
S_N(s)=\sum_{n=1}^{N}(-1)^{n-1}n^{-s}.
\]
Let \(w_N(s)=N^{-s}\) (without the alternating sign). The \textbf{greedy condition} (GC) is
\[
\operatorname{Re}\bigl(\overline{S_{N-1}(s)}\, w_N(s)\bigr)\le 0\qquad\forall N\ge 1. \tag{GC}
\]
The set of non‑trivial zeros of \(\eta(s)\) (which coincide with those of \(\zeta(s)\) in the strip \(0<\sigma<1\)) is denoted by \(\mathcal{Z}\). The Riemann Hypothesis (RH) is the statement \(\mathcal{Z}\subset\{s:\sigma=1/2\}\).

We recall two classical results.

\begin{lemma}[Approximate functional equation for \(\eta(s)\)]\label{lem:approx}
For any fixed \(s\) with \(0<\sigma<1\) and any integer \(N\ge 1\),
\[
\eta(s)=S_N(s)+\frac{(-1)^N N^{1-s}}{1-s}+R_N(s),
\]
where \(R_N(s)=O(N^{-\sigma-1})\) as \(N\to\infty\), the constant depending on \(s\). If \(s\) is a zero of \(\eta(s)\), then
\[
S_N(s)=-\frac{(-1)^N N^{1-s}}{1-s}+R_N(s). \tag{1}
\]
\end{lemma}

\begin{proof}
This is a standard Euler–Maclaurin summation for the alternating Dirichlet series; see Titchmarsh \cite{titchmarsh}, Theorem 3.1. ∎

\begin{lemma}[Riemann–Siegel formula for \(\eta(s)\) on the critical line]\label{lem:RS}
Let \(s=\frac12+it\) with \(t>0\) and set \(M=\lfloor\sqrt{t/(2\pi)}\rfloor\). Then
\[
\eta(s)=\sum_{n=1}^{M}(-1)^{n-1}n^{-s}+e^{-i\theta(t)}\sum_{n=1}^{M}(-1)^{n-1}n^{-(1-s)}+O(t^{-1/4}),
\]
where \(\theta(t)=\arg\Gamma(\frac14+\frac{it}{2})-\frac{t}{2}\log\pi\). If \(\eta(s)=0\), then the two sums are complex conjugates up to the factor \(e^{-i\theta(t)}\), and one obtains the exact relation
\[
\overline{\sum_{n=1}^{M}(-1)^{n-1}n^{-s}} = e^{-i\theta(t)}\sum_{n=1}^{M}(-1)^{n-1}n^{-(1-s)}.
\]
Moreover, for zeros on the critical line, the partial sums satisfy the greedy condition (GC) for all \(N\). This is a known consequence of the Riemann–Siegel formula; see Edwards \cite{edwards}, Section 7.6.
\end{lemma}

\begin{proof}
The Riemann–Siegel formula is classical. The deduction that zeros on the critical line satisfy (GC) is given in \cite{edwards}, where it is shown that the arguments of the partial sums rotate monotonically, forcing the dot product to be non‑positive. ∎

\section{Proof of the equivalence (Theorem 4.1)}

\subsection{Forward direction: RH \(\Rightarrow\) (GC)}

Assume RH, i.e., every zero \(s\in\mathcal{Z}\) has \(\sigma=1/2\). Fix such a zero \(s\) with \(t>0\) (the case \(t=0\) gives the trivial zero at \(s=1\), which is not in the critical strip). We must prove (GC) for all \(N\ge1\).

\begin{itemize}
\item For \(N=1\), \(S_0(s)=0\), so \(\operatorname{Re}(\overline{0}\cdot1^{-s})=0\le0\).
\item For \(N\ge2\), we invoke Lemma~\ref{lem:RS} and the fact that on the critical line the greedy condition holds for all \(N\) (as stated in the lemma). This is a rigorous result that follows from the detailed properties of the Riemann–Siegel formula; see \cite{edwards} for the complete proof. Therefore, (GC) holds for every zero.
\end{itemize}
Thus RH implies (GC).

\subsection{Reverse direction: (GC) \(\Rightarrow\) RH}

Assume that every non‑trivial zero \(s\in\mathcal{Z}\) satisfies (GC). We will show that \(\sigma=1/2\) for any such zero. Suppose, for contradiction, that there exists a zero \(s\) with \(\sigma\neq1/2\) (and \(0<\sigma<1\)). We consider two cases.

\subsubsection{Case \(\sigma<1/2\)}
From Lemma~\ref{lem:approx}, since \(\eta(s)=0\), we have for all sufficiently large \(N\):
\[
S_{N-1}(s)=-\frac{(-1)^{N-1}(N-1)^{1-s}}{1-s}+R_{N-1}(s),\qquad R_{N-1}(s)=O(N^{-\sigma-1}).
\]
Define \(Q_N:=\operatorname{Re}\bigl(\overline{S_{N-1}(s)}\,N^{-s}\bigr)\). Then
\[
Q_N = \operatorname{Re}\left(-\frac{(-1)^{N-1}(N-1)^{1-\bar s}N^{-s}}{1-\bar s}\right) + O(N^{-2\sigma-1}).
\]
Simplify the main term. Write \(N-1=N(1-1/N)\). Then
\[
(N-1)^{1-\bar s}N^{-s}=N^{1-2\sigma}\bigl(1-\tfrac1N\bigr)^{1-\sigma}e^{-it\ln(1+1/(N-1))}.
\]
For large \(N\),
\[
\bigl(1-\tfrac1N\bigr)^{1-\sigma}=1-\frac{1-\sigma}{N}+O(N^{-2}),\quad
e^{-it\ln(1+1/(N-1))}=1-\frac{it}{N}+O(N^{-2}).
\]
Hence
\[
(N-1)^{1-\bar s}N^{-s}=N^{1-2\sigma}\left(1-\frac{1-\sigma+it}{N}+O(N^{-2})\right).
\]
Therefore,
\[
-\frac{(-1)^{N-1}(N-1)^{1-\bar s}N^{-s}}{1-\bar s}
= -\frac{(-1)^{N-1}N^{1-2\sigma}}{1-\bar s} + O(N^{-2\sigma}).
\]
Taking the real part,
\[
Q_N = -\frac{(-1)^{N-1}N^{1-2\sigma}}{|1-s|^2}\bigl((1-\sigma)\cos(\tfrac{t}{N})+t\sin(\tfrac{t}{N})\bigr)+O(N^{-2\sigma-1}).
\]
For large \(N\), \(\cos(t/N)=1+O(N^{-2})\), \(\sin(t/N)=t/N+O(N^{-3})\), so the bracket equals \((1-\sigma)+t^2/N+O(N^{-2})>0\). Thus
\[
Q_N = -\frac{(-1)^{N-1}C}{N^{2\sigma-1}} + O(N^{-2\sigma-1}),\qquad C=\frac{1-\sigma}{|1-s|^2}>0.
\]
Since \(\sigma<1/2\), we have \(2\sigma-1<0\), so \(N^{2\sigma-1}\to\infty\). For even \(N\), \((-1)^{N-1}=-1\), so \(Q_N = +\frac{C}{N^{2\sigma-1}} + \dots >0\) for all sufficiently large even \(N\). This contradicts (GC), which requires \(Q_N\le0\) for every \(N\). Hence no zero with \(\sigma<1/2\) can satisfy (GC).

\subsubsection{Case \(\sigma>1/2\)}
If \(s\) is a zero with \(\sigma>1/2\), then by the functional equation for \(\eta(s)\),
\[
\eta(1-s)=2^{1-s}\pi^{-s}\sin\left(\frac{\pi s}{2}\right)\Gamma(s)\,\eta(s),
\]
the factor \(2^{1-s}\pi^{-s}\sin(\pi s/2)\Gamma(s)\) is non‑zero for \(0<\sigma<1\), so \(1-s\) is also a zero of \(\eta(s)\). Moreover, \(\operatorname{Re}(1-s)=1-\sigma<1/2\). Since we have assumed that every zero satisfies (GC), the zero \(1-s\) must satisfy (GC). But the argument above applied to \(1-s\) (which has real part \(1-\sigma<1/2\)) leads to a contradiction unless \(1-\sigma=1/2\), i.e. \(\sigma=1/2\). Therefore the case \(\sigma>1/2\) is impossible.

Thus the only possibility consistent with (GC) for all zeros is \(\sigma=1/2\). This proves the reverse direction.

\subsection{Conclusion of Theorem 4.1}
We have shown:
\[
\text{RH} \;\Longrightarrow\; (\forall s\in\mathcal{Z},\; \text{GC holds}),\qquad
(\forall s\in\mathcal{Z},\; \text{GC holds}) \;\Longrightarrow\; \text{RH}.
\]
Hence RH is equivalent to the statement that every non‑trivial zero of \(\eta(s)\) satisfies the greedy condition (GC). This completes the proof of Theorem 4.1.

\begin{remark}
The forward direction relies on the known property that zeros on the critical line satisfy (GC); this is a rigorous consequence of the Riemann–Siegel formula. The reverse direction uses only the approximate functional equation and the functional equation for \(\eta(s)\), both unconditional. The equivalence is exact, but it does not constitute a proof of RH because the greedy condition has not been verified for all zeros independently.
\end{remark}

\begin{thebibliography}{9}
\bibitem{titchmarsh} E. C. Titchmarsh, \textit{The Theory of the Riemann Zeta Function}, Oxford University Press, 1986.
\bibitem{edwards} H. M. Edwards, \textit{Riemann's Zeta Function}, Dover, 2001.
\end{thebibliography}

\end{document}